\subsection{A staged Morita obstruction for \texorpdfstring{$R_{7,4,2591}$}{R7-4-2591}}
\label{sec:r2591-staged-morita}

\begin{theorem}\label{thm:r2591-staged-morita}
The fusion ring $R_{7,4,2591}$ admits no categorification over $\C$.
No pivotality or unitarity assumption is required.
\end{theorem}

\begin{proof}
Put $t=2+\sqrt5$, so $t^2=1+4t$.  In a hypothetical
categorification $\mathcal C$, the invertibles form $G\cong V_4$;
the other simples are $X,Y,Z$, of dimensions $2,2t,2t$, with
$Y^*=Z$.  Every element of $G$ fixes each of these three objects on
both sides, and
\[
 X^2=\sum_{g\in G}g,\qquad XY=YX=2Z,\qquad
 XZ=ZX=2Y,\qquad Y^2=2X+4Y+4Z.
\]
Proposition~\ref{prop:twisted-v4} supplies a Morita equivalent category
$\mathcal E$ with eight invertibles $K$ and eight dimension-$t$
simples forming a free $K$-orbit. It is the dual of the nontrivially
twisted regular algebra $P$ on $G$.

The inverse algebra $Q$ is the regular object on a subgroup
$L\cong V_4$ of $K$, with possibly twisted multiplication.  Indeed,
the internal end of the distinguished right $P$-module $P$ contains
each of the four character-twist bimodules on pointed support once
and no other simple: its multiplicity on a $P$-bimodule $V$ is
$\dim\operatorname{Hom}_{\mathcal C}(\mathbf1,V)$.
Consequently $\mathcal C\simeq{}_Q\mathcal E_Q$.  The existence of
$Q$ trivializes the associator on $L$, but no triviality outside $L$
is assumed.

\medskip\noindent\emph{A central fixed involution.}
Choose a simple $\rho$ in the free orbit and write
\[
 \rho g=\phi(g)\rho,\qquad \rho^2=a+N\rho.
\]
Associativity gives $\phi\in\operatorname{Aut}(K)$,
$\phi^2=\operatorname{Ad}(a)$ and $\phi(a)=a$; also
$\rho^*=a^{-1}\rho$.  The subgroup $L$ has index two and is normal.
We claim $\phi(L)=L$.  Otherwise $L\phi(L)=K$, so the noninvertible
$L$-by-$L$ action has a single orbit of size eight and an order-two
stabilizer with injective projections into $L$.  Its two projective
characters give $Y,Z$, since $H^2(C_2,\C^\times)=0$.
The character twists of $Q$ are precisely the four original
invertibles.  Their characters restrict surjectively to this
stabilizer, so some such invertible interchanges $Y$ and $Z$,
contradicting their individual fixedness.  This argument is unchanged
by the multiplication cocycle of $Q$.

There are therefore two noninvertible support orbits, $L\rho$ and
$(K\setminus L)\rho$, each supporting exactly one of $Y,Z$.
Their mutual duality forces $a\notin L$.  The only groups of order
eight containing $V_4$ are $C_4\times C_2$, $C_2^3$ and $D_8$.
The last is impossible: conjugation by any element outside an
index-two $V_4$ in $D_8$ transposes its two reflections, whereas
$(\phi|_L)^2$ cannot be a transposition in
$\operatorname{Aut}(V_4)=S_3$.  Thus $K$ is abelian and $\phi^2=1$.
The restriction $\phi|_L$ fixes a nonidentity element $\delta$.
It is central in $K$, and $\phi(\delta)=\delta$.

\medskip\noindent\emph{Factoring the inverse algebra.}
The summand $A=\mathbf1+\delta$ is an actual subalgebra of $Q$;
its multiplication is the restriction of the twisted regular
multiplication.  It is connected and separable.  In
$\mathcal E'={}_A\mathcal E_A$, all pointed and noninvertible
support orbits have size two and stabilizer $C_2$, because $\delta$
is central and $\phi$-fixed.  Each orbit splits into two projective
characters.  Thus $\mathcal E'$ again has eight invertibles $K'$
and eight noninvertibles of dimension $t$.  The same dimension
comparison shows that $K'$ acts freely and transitively on them.

Regard $Q$ as an algebra $B$ in $\mathcal E'$.  Its multiplication
factors through $Q\otimes_A Q$, its unit is $A\longrightarrow Q$,
and separability descends to the relative tensor product.  It has
$\FPdim_{\mathcal E'}(B)=4/2=2$.  Moreover evaluation at the unit
injects
\[
 \operatorname{Hom}_{A\text{-}A}(A,Q)
 \longrightarrow \operatorname{Hom}_{\mathcal E}(\mathbf1,Q),
\]
whose target is one-dimensional; the inclusion $A\to Q$ supplies
a nonzero map.  Thus $B$ is connected.  Since $t>2$, it must be
$B=\mathbf1+h$ for an invertible involution $h$.  This algebra is
already supplied by $Q$, rather than inferred from a based-ring
subgroup.  Restriction and extension of the two $Q$ actions give
the canonical relative-Morita identification
\[
 {}_B({}_A\mathcal E_A)_B\simeq{}_Q\mathcal E_Q\simeq\mathcal C.
\]

\medskip\noindent\emph{The forced dihedral support permutation.}
Put $H=\langle h\rangle\subset K'$.  If $H$ were normal, all four
pointed double cosets would have order-two stabilizer and would
yield eight invertibles.  This contradicts the integral part of
$\mathcal C$, whose dimensions are $1,1,1,1,2$.
Therefore $K'=D_8=\langle r,s\mid r^4=s^2=1,\ srs=r^{-1}\rangle$
and we may take $H=\langle s\rangle$.
For a noninvertible $\sigma$ write $\sigma g=\psi(g)\sigma$.
The stabilizer of the support through $g\sigma$ is
$H\cap g\psi(H)g^{-1}$.  If $H$ and $\psi(H)$ were conjugate,
some support orbit would have size two and produce simples of
dimension $t$, which do not occur in $\mathcal C$.
They lie in opposite reflection classes; changing generators gives
$\psi(H)=\langle rs\rangle$.

The two noninvertible support orbits are
\[
 U\sigma,\quad r^2U\sigma,
 \qquad U=\{1,s,rs,r^{-1}\}.
\]
The pointed coset $Hr^2$ has two invertible bimodules, and tensoring
either of them multiplies noninvertible support by the central
rotation $r^2$.  It therefore interchanges the two supports, hence
interchanges $Y$ and $Z$.  This contradicts their individual
fixedness under every original invertible and proves the theorem.
\end{proof}

\begin{remark}
The source-bound certificate
\texttt{verify\_R2591\_staged\_morita.py} replays the twisted-block
calculation, the source products and duality, all five groups of
order eight, and all sixteen opposite-reflection configurations.
The compatible inverse-group systems comprise eight for
$C_4\times C_2$ and $112$ for $C_2^3$, with none for $D_8$.
These finite checks verify the algebraic data used in the proof;
the actual-algebra and relative-Morita arguments above are
categorical and have been independently audited.
\end{remark}
